\chapter{Diskussion}
\label{chap:discussion}
Der folgende Abschnitt soll die Ergebnisse rekapitulieren, diskutieren und ihre Limitationen aufzeigen. \\
Im ersten Vorexperiment sollten die Detektoren auf einer Stichprobe von 1.000 Posts getestet werden. Hier zeigt sich, dass die Leistung der Detektoren auf der Stichprobe stark variiert. Am besten schneidet hierbei jedoch Binoculars ab, das über alle beobachteten Metriken die höchsten Werte erreicht und menschliche als auch KI-generierte Texte mit einer hohen Zuverlässigkeit klassifiziert. RADAR und der ChatGPT Detector RoBERTa zeigen entgegengesetzte Muster. Ersterer weist einen hohen Recall auf und erkennt somit viele KI-Texte korrekt. Gleichzeitig werden viele menschliche Texte fälschlich als KI-generiert klassifiziert. Der ChatGPT Detector RoBERTa erkennt die überwiegende Mehrheit der menschlichen Texte als solche. KI-Texte werden von diesem Detektor aber häufig nicht korrekt erkannt und als menschlich klassifiziert. Diese Ergebnisse zeigen, dass die Analyse einer einzelnen Evaluationsmetrik zu kurz greift, um die gesamte Erkennungsleistung zu beurteilen. Auffallend ist, dass der RoBERTa Base OpenAI Detector trotz seiner vergleichsweise veralteten und eingeschränkten Trainingsdaten hinsichtlich der betrachteten Metriken unter den experimentellen Bedingungen eine bessere Erkennungsleistung zeigt als der jeweils schlechteste Einzeldetektor. In Bezug auf die Zahl der False Negatives zeigen RADAR und Binoculars die niedrigsten Werte. Binoculars ist dabei allerdings besser in der Lage gleichzeitig auch die False Positive Rate niedrig zu halten. Somit ist Binoculars unter den im Vorexperiment untersuchten Bedingungen der zuverlässigste Detektor. Da hier allerdings nur ein Ausschnitt der Daten verwendet wurde und die Generierung ausschließlich mit LLaMA-3.1-8B-Instruct erfolgt ist, können die Ergebnisse nicht ohne Weiteres auf andere Texte und generierende Sprachmodelle übertragen werden.
Das Ensemble mit Mehrheitsentscheidung ist für die Metriken Accuracy, Precision und F1-Score leistungsstärker als alle Einzeldetektoren mit der Ausnahme von Binoculars.
Der deutlich geringere Recall des Ensembles im Vergleich zu Binoculars weist darauf hin, dass bei einem beträchtlichen Teil der KI-generierten Texte nicht mindestens drei der vier Detektoren gleichzeitig für eine KI-Klassifikation stimmen. Es könnte sein, dass die hohe Zahl der False Negatives des ChatGPT Detector RoBERTa hierzu beiträgt. Wenn das Ensemble einen Text als KI-generiert erkennt, stimmt diese Entscheidung im Großteil der Fälle. Allerdings werden über 40 \% der KI-generierten Texte nicht als solche erkannt. Auf Basis dieser Stichprobe und den untersuchten Bedingungen kann das hier zugrundeliegende Ensemble mit einer \(\frac{3}{4}\)-Mehrheitsentscheidung nicht die Leistung des starken Einzeldetektors Binoculars übertreffen. Die hier verwendete Kombination mehrerer Detektoren zu einem Ensemble mit \(\frac{3}{4}\)-Mehrheitsentscheidung kann also die Schwächen der Einzeldetektoren nicht ausreichend ausgleichen, auch nicht bei der Verwendung von Binoculars als starkem Detektor. Im hier durchgeführten Vorexperiment gehen die Detektoren alle gleichermaßen in die Mehrheitsentscheidung ein. Alternativ könnte man die performanteren Detektoren höher gewichten und so deren Einfluss auf das Mehrheitsvotum stärken.
Unter diesen Bedingungen lässt sich die Hypothese \hyperlink{hp:2}{H\textsubscript{2}} vorläufig nicht bestätigen. Allerdings müssen diese Ergebnisse vor dem Hintergrund der Beschränkung auf Posts aus nur einem einzelnen Monat und der ausschließlichen Verwendung von LLaMA-3.1-8B-Instruct zur Textgenerierung gesehen werden. \\
Im zweiten Vorexperiment sollten KI-Verdächtigungen in den Kommentaren quantifiziert werden. Dies sollte anhand einer Schlagwortsuche mit Hilfe von Regular Expressions erfolgen. Im Ergebnis zeigte sich anhand der verwendeten Regex-Muster im gesamten Untersuchungszeitraum zwischen Januar 2022 und Dezember 2024, dass nur ein sehr kleiner Teil der Kommentare als KI-Verdächtigung identifiziert wird. Prinzipiell lässt sich zwar ein leichter Anstieg über die Zeit hinweg beobachten, allerdings weist der Verlauf starke Schwankungen auf. Dieser leichte Anstieg könnte in Zusammenhang mit der zunehmenden Verbreitung von LLMs stehen. Des Weiteren können sich die Formulierungen und Begriffe in diesem Kontext verändert haben. Beispielhaft dafür ist der Begriff „ai slop“, dessen Verwendung im Laufe der Jahre in englischsprachigen sozialen Medien deutlich zugenommen hat und ältere ähnlich abwertende Begriffe in diesem Kontext abgelöst hat \autocite{miklian2026s}. In diesem Sinne könnte ein Teil des Anstiegs damit erklärt werden, dass Begriffe wie „ChatGPT“ und „ai-generated“ erst mit der zunehmenden Prominenz von ChatGPT und generativen KI-Modellen im Allgemeinen vermehrt in den Sprachgebrauch aufgenommen wurden. Generell lässt sich in bestimmten Subreddits eine ablehnende Haltung gegenüber KI-generierten Inhalten erkennen, was dazu führen könnte, dass häufiger KI-Verdächtigungen ausgesprochen werden. Dies könnte insbesondere der Fall bei Subreddits sein, die auf menschlich produzierte Inhalte Wert legen und diese auch in ihren Regeln festlegen, wie „r/changemyview“ oder „r/WritingPrompts“ \autocite{lloyd2025there}.
Insgesamt sind die absoluten Häufigkeiten von KI-Verdächtigungen sehr gering und erschweren damit die Interpretation der zeitlichen Entwicklung. Die vergleichsweise großen Unterschiede in den relativen Werten können bereits durch geringe Schwankungen in den absoluten Zahlen der identifizierten KI-Verdächtigungen ausgelöst werden. In Anbetracht dessen muss auch die Interpretation des leichten Anstiegs mit Bedacht erfolgen. Ein Grund für die geringe absolute Anzahl an KI-Verdächtigungen im vorliegenden Datensatz, könnte die Auswahl der Regex-Muster sein. Diese berücksichtigen ausschließlich deutsch- und englischsprachige Kommentare und decken nicht alle in den Reddit-Kommentaren vorhandenen Sprachen ab. Es könnte also sein, dass die Zahl der KI-Verdächtigungen durchaus höher liegt, wenn die Regex-Muster auch um andere Sprachen ergänzt werden. Außerdem könnte es sein, dass die konkret untersuchten Formulierungen deutlich von den gängigen auf Reddit verwendeten Phrasen abweichen. Kommentare können zwar die untersuchten Phrasen enthalten, aber dennoch keine KI-Verdächtigung als solche gegenüber einem Post oder Kommentar äußern. Dies kann der Fall sein, wenn der Verfasser des Kommentars eine Phrase nutzt, die aber im Kontext einer allgemeinen Diskussion zum Thema KI verwendet wird, ohne einen konkreten Verdacht zu äußern. Zudem können besonders kurze Muster fälschlicherweise Treffer in anderen Sprachen als Englisch oder Deutsch erzeugen, wenn diese dort eine andere Bedeutung haben (Beispiel: „ai post“ auf Italienisch). Die verwendeten Regular Expressions können somit sowohl zu False Positives als auch False Negatives führen. Zusammenfassend kann gesagt werden, dass die hier festgestellte Anzahl an Treffern auf Grundlage von Regular Expressions nicht zwangsläufig der Anzahl an tatsächlich geäußerten Verdächtigungen entspricht. Das Vorexperiment stellt lediglich einen Versuch der Abschätzung dar und erfasst zunächst nur sprachliche Hinweise auf KI-Verdächtigungen. Auf Basis der Ergebnisse des zweiten Vorexperiments stellt sich auch die Frage, wie sich KI-Verdächtigungen in natürlicher Sprache in Form von Kommentaren identifizieren lassen. Eine einfache Schlagwortsuche, wie sie hier verwendet worden ist, könnte nicht ausreichend sein, um eine zuverlässige Erkennung zu ermöglichen. Es wäre denkbar, dass eine Zero-Shot-Klassifikation über Natural Language Inference (NLI)-Ansätze, die Verwendung von LLMs oder ein Fine-Tuning bestehender Klassifikatoren eine sinnvolle Ergänzung darstellen könnten. Diese sind eventuell besser in der Lage bestimmte potenzielle Phrasen der Verdächtigung im Kontext zu sehen. Zu erwähnen ist auch, dass sich diese Ergebnisse nicht auf die Grundgesamtheit aller auf Reddit veröffentlichten Kommentare übertragen lassen. Es werden aus bestimmten Monaten nur Stichproben betrachtet und manche Monate enthalten ohnehin nur 100.000 Kommentare. \\
Das erste Vorexperiment diente als Maßgabe für das weitere Vorgehen in dieser Arbeit. Da vor allem Binoculars auf der untersuchten Stichprobe gute Ergebnisse erzielt hat und auch die anderen Modelle in gewissen Bereichen ihre Stärken haben, sollen alle Detektoren im Hauptexperiment auf einer größeren Stichprobe evaluiert werden. Hierzu wurden zunächst verschiedene Zeitpunkte betrachtet. Diese Zeitpunkte sollen sich an Meilensteinen in der Entwicklung und der Verbreitung von LLMs orientieren. Für alle Zeitpunkte wird zur Generierung immer LLaMA-3.1-8B-Instruct verwendet. Es wird also nicht unmittelbar die technische Weiterentwicklung der LLMs über die Zeit hinweg abgebildet. Vielmehr wird untersucht, ob ein Zusammenhang zwischen dem Veröffentlichungszeitraum der zugrunde liegenden Posts bzw. deren möglicher Veränderung über die Zeit und der Leistung der Detektoren besteht. Dabei diente der Juni 2022 als Referenzzeitpunkt vor der Veröffentlichung von ChatGPT. Es kann davon ausgegangen werden, dass zu diesem Zeitpunkt die Verbreitung von LLMs sehr beschränkt war. Die darauffolgenden Zeitpunkte wurden so gewählt, dass sie die großen Weiterentwicklungsschritte von LLMs abbilden. Aus diesem Grund wurden der Dezember 2022 (Veröffentlichung ChatGPT), der März 2023 (Veröffentlichung GPT-4), der Juli 2023 (Veröffentlichung LLaMA 2), der Januar 2024 und der Juli 2024 (Veröffentlichung LLaMA-3.1) als weitere Zeitpunkte untersucht. Aus jedem dieser Monate wurde eine Stichprobe von 3.000 Posts gezogen, auf deren Basis dann ebenso wie in Vorexperiment 1 mithilfe von LLaMA-3.1-8B-Instruct künstliche Posts generiert wurden. \\
Zunächst sollte vermutet werden, dass so gut wie alle menschlichen Posts in der Stichprobe aus dem Juni 2022 von den Detektoren auch als solche klassifiziert werden. Allerdings haben bereits zu diesem Zeitpunkt alle Detektoren in mehr oder weniger großem Ausmaß menschlichen Text als KI erkannt. Besonders auffällig ist in diesem Kontext die False Positive Rate von RADAR. Dieser Detektor erkennt in der Stichprobe aus dem Juni 2022 67 \% der von Menschen verfassten Posts fälschlicherweise als KI-generiert. Auch über alle anderen untersuchten Zeiträume weist RADAR die mit Abstand höchste FPR auf. Hier wird dieselbe Problematik verdeutlicht, die bereits im ersten Vorexperiment beobachtet werden konnte. RADAR hat nach Binoculars den höchsten Recall und klassifiziert gleichzeitig auch einen erheblichen Teil der menschlichen Texte fälschlicherweise als KI-generiert. Dies legt nahe, dass RADAR unter den hier verwendeten experimentellen Bedingungen bestimmte Merkmale von menschlichen Reddit-Posts als Indikatoren für KI-generiertem Text interpretiert. In weniger ausgeprägter Form könnte dies auch bei den anderen Detektoren der Fall sein. Die FPR des ChatGPT Detector RoBERTa und von Binoculars bewegen sich über die betrachteten Zeitpunkte hinweg auf einem ähnlich niedrigen Niveau. Beide Detektoren erkennen also menschlichen Text am zuverlässigsten auch als solchen. Gleichzeitig gibt es aber einen deutlichen Unterschied im Recall. Der ChatGPT Detector RoBERTa erkennt nur rund 20 \% der KI-generierten Texte auch als solche, während Binoculars einen Recall von deutlich über 90 \% erreicht. Die hohe Diskrepanz zwischen der FPR und dem Recall des ChatGPT Detector RoBERTa deutet darauf hin, dass der Detektor in seiner Klassifikation eher zurückhaltend ist. Nur wenige menschliche Texte werden fälschlich als KI-generiert erkannt und gleichzeitig werden nur wenige der tatsächlich KI-generierten Texte als solche klassifiziert. Auch hier zeigt sich, dass eine Metrik alleine nicht ausreicht, um die Leistung eines Detektors hinreichend zu beurteilen.\\
Über die untersuchten Zeiträume hinweg lässt sich nur eine sehr schwache Veränderung in der Erkennungsleistung beobachten. Die Metriken verhalten sich für alle evaluierten Detektoren überwiegend stabil. Außerdem sind in der untersuchten Stichprobe keine wesentlichen Unterschiede in der Rangfolge der Detektoren erkennbar. Es zeigt sich unter den gegebenen experimentellen Bedingungen und der gewählten Stichprobe keine auffällige Leistungssteigerung oder -abfall der Detektoren über die Zeit hinweg. Zwar lassen sich für einzelne Detektoren kleine Unterschiede hinsichtlich der FPR zwischen den Stichproben aus dem Juni und dem Dezember 2022 erkennen, aber daraus kann nicht ohne Weiteres ein Effekt der Veröffentlichung von ChatGPT geschlossen werden. Dazu müssten zusätzliche Zeitpunkte vor und nach Dezember 2022 betrachtet werden. Es kann zudem nicht ausgeschlossen werden, dass diese Fluktuationen auf unterschiedliche Subreddits, Textlängen und allgemein Stichprobenzusammensetzungen zurückzuführen sind. Hierzu müsste eine Kontrolle dieser Einflussfaktoren erfolgen. Auch hinsichtlich der Unterschiede im Recall der Detektoren zwischen den Stichproben aus Januar und Juli 2024 lässt sich nicht schlussfolgern, dass die Leistung der Detektoren über die Zeit hinweg abnimmt. In der Stichprobe aus Juli 2024 steigen die Werte wieder an und der vorhergehende Rückgang kann ebenfalls durch andere Eigenschaften des untersuchten Textes wie Textlänge und Subreddit beeinflusst worden sein.
Insgesamt zeigt sich Binoculars auch auf der größeren Stichprobe über alle beobachteten Monate und verwendeten Metriken hinweg als der zuverlässigste Detektor. Die Unterschiede zwischen den Detektoren scheinen stärker zu sein als die Differenzen zwischen den hier untersuchten Zeitpunkten. Es lässt sich also unter den gegebenen experimentellen Bedingungen kein eindeutiger Zusammenhang zwischen dem Veröffentlichungszeitraum der Reddit-Posts und der Erkennungsleistung der Detektoren feststellen. Zu berücksichtigen ist dabei, wie oben geschildert, dass für alle beobachteten Zeitpunkte dasselbe Modell, LLaMA-3.1-8B-Instruct, zur Generierung verwendet wurde. Die Ergebnisse zeigen daher nicht, wie sich die Leistung der Detektoren im Vergleich zu den jeweils tatsächlich verfügbaren LLMs entwickelt hat. Vielmehr zeigen sie, inwiefern die Erkennungsleistung variiert, wenn die menschlichen Reddit-Posts aus unterschiedlichen Veröffentlichungszeiträumen stammen. \\
Im zweiten Teil des Hauptexperiments sollte der Einfluss der Textlänge auf die Detektionsleistung untersucht werden. Dazu wurden zunächst verschiedene Klassen an Textlängen gebildet (0-150, 150-300, 300-450, 450-600, 600-750, 750-900 Zeichen). Anschließend wurde für jede Textlängenklasse eine Stichprobe von jeweils 3.000 realen und KI-generierten Posts gebildet. Die Posts stammen aus denselben Monaten, die bereits im ersten Teil des Hauptexperiments verwendet wurden. Grundsätzlich ist ein positiver Zusammenhang zwischen der Textlänge und der Erkennungsleistung erkennbar. Allerdings gibt es Unterschiede zwischen den beobachteten Detektoren und den untersuchten Metriken. Insgesamt zeigt sich, dass mit längeren Texten offenbar auch mehr Informationen geliefert werden, die von den verwendeten Detektoren sinnvoll zur Klassifikation genutzt werden können. Sind die Texte sehr kurz, haben die Detektoren nur wenige Merkmale zur Verfügung anhand derer sie menschlichen von KI-generiertem Text unterscheiden können. Je länger der Text wird, desto eher steigt die Zahl von potenziell relevanten Mustern in Vokabular, Syntax und statistischen Regelmäßigkeiten. Betrachtet man die Accuracy und F1-Score fällt auf, dass vor allem Binoculars mit einer zunehmenden Textlänge auch deutlich bessere Ergebnisse liefert. Dies könnte dadurch erklärt werden, dass dieser Detektor diese Muster in längerem Text besser erkennt und nutzen kann. Für kürzere Texte weist Binoculars allerdings eine ähnlich eingeschränkte Zuverlässigkeit auf wie die anderen untersuchten Detektoren. Binoculars zeigt interessanterweise auch eine sichtliche Steigerung der Werte für den Recall und die Precision. Somit werden bei längeren Texten erstens mehr KI-generierte Texte richtig klassifiziert und zweitens auch mehr KI-Texte insgesamt erkannt. Auch dies deutet darauf hin, dass die Zuverlässigkeit von Binoculars bei kurzen Reddit-Posts nur sehr limitiert ist und in diesen Fällen die Ergebnisse deshalb nur mit Vorsicht interpretiert werden sollten. Anders verhält sich RADAR hinsichtlich Recall und FPR. Ähnlich wie oben bei Betrachtung zu unterschiedlichen Zeitpunkten weist RADAR hier über alle betrachteten Textlängen hinweg den höchsten Recall auf. Dass RADAR bei kurzen Texten einen hohen Recall zeigt, hängt vermutlich nicht damit zusammen, dass der Detektor besonders gut in der Lage ist KI-generierten Text zu erkennen, sondern eher, dass RADAR generell häufig für das Label „KI-generiert“ stimmt. Dafür spricht auch die sehr hohe FPR bei kurzen Texten. Diese sinkt zwar mit zunehmender Textlänge sichtlich, verläuft aber auch bei längeren Posts deutlich oberhalb der FPR der anderen untersuchten Detektoren. Werden die Texte länger, scheint RADAR für sich genommen besser in der Lage zu sein zwischen menschlichen und KI-generierten Posts zu differenzieren. Der Kurvenverlauf für die Accuracy, die Precision und den F1-Score des RoBERTa Base OpenAI Detectors spricht für diesen Detektor dafür, dass der Zusammenhang zwischen der Textlänge und der Detektionsleistung nicht unbedingt monoton ist. Auch hier könnte nicht nur die Textlänge einen Einfluss haben, sondern daneben auch Komponenten wie Stichprobenzusammensetzung oder Subreddit-spezifische Eigenheiten eine Rolle spielen. Für das Ensemble gelten hinsichtlich der Korrelation zwischen Textlänge und Erkennungsleistung ähnliche Tendenzen. Mit längeren Texten geht zwar eine bessere Erkennungsleistung einher, aber diese überbietet nicht die stärksten Einzeldetektoren. Insgesamt profitieren zwar alle untersuchten Modelle von längeren Texten und damit mehr Informationen, aber die Schwächen einzelner Detektoren werden aufgrund der \(\frac{3}{4}\)-Mehrheitsentscheidung auch hier nicht hinreichend ausgeglichen. Außerdem lässt sich unter den hier betrachteten experimentellen Bedingungen kein Kausaleffekt der Textlänge auf die Erkennungsleistung attestieren. Dazu bedarf es der Kontrolle anderer möglicher Einflussfaktoren.
